\newpage
\section{Operaciones Bancarias Seguras}

\subsection{Atajos de latencia}

Durante las pruebas, nos dimos cuenta de que no siempre hace falta recurrir al 
modelo de lenguaje. Por ejemplo, cuando el usuario simplemente pregunta 
"¿cuál es mi saldo?", pasar la pregunta por el LLM añadía unos segundos de 
latencia innecesaria para devolver un dato directo.

Para optimizar esto, en el archivo \texttt{agent.py} incluimos la función 
\texttt{es\_consulta\_saldo}. Si el mensaje coincide con esta validación, 
el \emph{backend} llama directamente a la API \texttt{consultar\_saldo} y 
responde al momento, saltándose por completo al modelo.

El filtro funciona detectando palabras clave (como saldo o dinero disponible) 
y comprobando que el resto de la frase solo contenga términos de relleno autorizados 
en la lista \texttt{\_RELLENO\_SALDO}. Hemos sido muy restrictivos a propósito con estas 
palabras permitidas, para evitar que consultas que requieran información adicional 
sean identificadas incorrectamente como una consulta de saldo.

En este caso, los errores tienen distinto coste:
\begin{itemize}
    \item \textbf{Falso negativo:} si una consulta de saldo simple no pasa el filtro, 
    la petición llegará al LLM, que usará la herramienta correspondiente y responderá 
    correctamente, solo que tardará un poco más.

    \item \textbf{Falso positivo:} si confundimos una pregunta compleja con una simple, 
    el sistema fallará. Por ejemplo, si el usuario pregunta "¿cuál era mi saldo el mes 
    pasado?", el atajo le devolvería el saldo actual, dando una respuesta completamente 
    incorrecta.
\end{itemize}

Por este motivo, ante la más mínima duda (si hay palabras fuera de la lista de relleno), 
el filtro prefiere no aplicar el atajo y dejar que el modelo analice la intención completa 
de la frase de forma segura.

\subsection{Flujo transaccional de Bizum}

Mover dinero es la operación más crítica de la aplicación. Por seguridad, la ejecución de un 
Bizum nunca se delega libremente al modelo de lenguaje; el LLM o las expresiones regulares 
solo ayudan a entender la intención, pero la ejecución está controlada de principio a fin por 
el \emph{backend}. El proceso, ilustrado en la Figura~\ref{fig:bizum}, sigue estas etapas:

\begin{figure}[htbp]
\centering
\begin{tikzpicture}[
    font=\small,
    caja/.style={draw=#1, fill=#1!7, rounded corners=2pt, align=center,
                 minimum width=2.9cm, minimum height=0.95cm, inner sep=3pt},
    proceso/.style={caja=subsectionblue},
    exito/.style={draw=subsectionblue, fill=subsectionblue!20, rounded corners=2pt,
                  align=center, minimum width=2.9cm, minimum height=0.95cm,
                  inner sep=3pt, thick},
    rechazo/.style={caja=ugrred},
    aviso/.style={caja=gray},
    decision/.style={chamfered rectangle, chamfered rectangle xsep=4pt,
                     draw=gray!70!black, fill=gray!8, align=center,
                     minimum width=2.9cm, minimum height=0.9cm, inner sep=3pt},
    flecha/.style={-{Stealth[length=2mm]}, thick, draw=gray!70!black},
    et/.style={font=\scriptsize, text=gray!55!black, inner sep=2pt},
]

% --- Camino principal ------------------------------------------------------
\node[proceso] (ini) at (0,0)
      {Petición de Bizum\\[-1pt]\scriptsize regex del backend o \textit{tool call} del LLM};
\node[decision] (d1) at (0,-1.8)  {¿Indica el importe?};
\node[decision] (d2) at (0,-3.6)  {¿Entre 0,50 y 1.000 €?};
\node[decision] (d3) at (0,-5.4)  {¿Existe el contacto?};
\node[proceso]  (pe) at (0,-7.2)
      {Bizum pendiente\\[-1pt]\scriptsize evento \texttt{pedir\_pin}};
\node[decision] (d4) at (0,-9.0)  {¿PIN correcto?};
\node[proceso]  (tx) at (0,-10.8)
      {Transacción \texttt{BEGIN IMMEDIATE}\\[-1pt]\scriptsize límite diario de 500 € y saldo};
\node[exito]    (ok) at (0,-12.6)
      {Descuenta el saldo\\[-1pt]\scriptsize eventos \texttt{saldo} y ticket};

\draw[flecha] (ini) -- (d1);
\draw[flecha] (d1) -- node[et, right] {sí} (d2);
\draw[flecha] (d2) -- node[et, right] {sí} (d3);
\draw[flecha] (d3) -- node[et, right] {sí} (pe);
\draw[flecha] (pe) -- (d4);
\draw[flecha] (d4) -- node[et, right] {sí} (tx);
\draw[flecha] (tx) -- node[et, right] {cabe} (ok);

% --- Salidas a la derecha --------------------------------------------------
\node[aviso]   (r1) at (5.6,-1.8)  {Pregunta el importe\\[-1pt]\scriptsize no deja nada pendiente};
\node[rechazo] (r2) at (5.6,-3.6)  {Rechaza\\[-1pt]\scriptsize sin pedir el PIN};
\node[rechazo] (r3) at (5.6,-5.4)  {Avisa y termina};
\node[aviso]   (r4) at (5.6,-9.0)  {Resta un intento\\[-1pt]\scriptsize 3 como máximo};
\node[rechazo] (c4) at (5.6,-10.8) {Cancela el envío};

\draw[flecha] (d1) -- node[et, above] {no} (r1);
\draw[flecha] (d2) -- node[et, above] {no} (r2);
\draw[flecha] (d3) -- node[et, above] {ninguno} (r3);
\draw[flecha] (d4) -- node[et, above] {no} (r4);
\draw[flecha] (r4.north) |- node[et, right, pos=0.25] {quedan intentos} (pe.east);
\draw[flecha] (r4) -- node[et, right] {sin intentos} (c4);

% --- Salidas a la izquierda ------------------------------------------------
\node[decision] (sg) at (-5.6,-5.4)  {¿Querías decir…?};
\node[rechazo]  (c3) at (-5.6,-3.6)  {Cancela};
\node[rechazo]  (r5) at (-5.6,-10.8) {Rechaza\\[-1pt]\scriptsize la BD no cambia};

\draw[flecha] (d3) -- node[et, above] {parecido} (sg);
\draw[flecha] (sg) -- node[et, right] {no} (c3);
\draw[flecha] (sg.south) |- node[et, right, pos=0.25] {sí} (pe.west);
\draw[flecha] (tx) -- node[et, above] {no cabe} (r5);

% --- Nota sobre el PIN -----------------------------------------------------
\node[font=\scriptsize, text=ugrred, align=center, draw=ugrred!60, dashed,
      rounded corners=2pt, inner sep=4pt] at (-5.6,-8.6)
      {El PIN viaja por el evento\\\texttt{auth\_bizum}: el modelo\\nunca lo recibe};

\end{tikzpicture}
\caption{Flujo de un envío de Bizum. Todas las comprobaciones se hacen en el
backend y en este orden: el importe se valida antes de pedir el PIN, y el
límite diario y el saldo se comprueban dentro de la misma transacción que los
descuenta.}
\label{fig:bizum}
\end{figure}

\begin{enumerate}

    \item \textbf{Extracción y validación inicial:} la función
    \texttt{extraer\_peticion\_bizum} identifica en la petición del usuario los
    datos necesarios para realizar el envío, como el destinatario, el importe
    y el concepto. La extracción permite reconocer diferentes formas de
    expresar una misma operación. Antes de continuar con el proceso, el
    importe se comprueba para verificar que se encuentra dentro de los límites
    establecidos, entre 0,50€ y 1000€.

    \item \textbf{Resolución del contacto:} una vez obtenido el destinatario, llamamos a
    la función \texttt{buscar\_contacto\_bizum}, que comprueba los contactos
    disponibles. En caso de que el nombre proporcionado no coincida
    exactamente, se utilizan coincidencias parciales mediante
    \texttt{get\_close\_matches} para identificar posibles contactos. Si existe
    una coincidencia, esta se muestra al usuario antes de continuar con la
    operación.

    \item \textbf{Autenticación desde el frontend:} cuando los datos de la
    operación son válidos, el estado de la conversación pasa a ``Bizum
    pendiente'' y el \emph{backend} envía mediante WebSocket la señal
    \texttt{pedir\_pin}. El \emph{frontend} muestra entonces un teclado numérico para
    introducir el PIN. De esta forma, el PIN se introduce fuera del campo de
    texto de la conversación y no forma parte del historial que recibe el
    modelo de lenguaje.

    \item \textbf{Validación y ejecución de la operación:} una vez validado el
    PIN, el \emph{backend} ejecuta la función \texttt{api\_enviar\_bizum}. La
    modificación del saldo se realiza dentro de una transacción SQLite
    iniciada mediante \texttt{BEGIN IMMEDIATE}, utilizando
    \texttt{transaccion\_escritura}. Esta instrucción establece el bloqueo
    necesario para evitar que dos operaciones de escritura se ejecuten
    simultáneamente sobre los mismos datos. De esta forma, la comprobación del
    saldo y su posterior actualización se realizan dentro de la misma
    transacción.

\end{enumerate}

Este diseño mantiene separadas la interpretación de la petición y la ejecución
de la operación bancaria. El modelo de lenguaje puede identificar la intención
del usuario y proporcionar los datos necesarios, pero las validaciones y la
modificación de los datos se realizan mediante la lógica del backend.